\chapter{Troubleshooting and FAQ}\label{ch:faq}

\newcommand{\faq}[1]{\needspace{4\baselineskip}\subsection*{\faQuestionCircle[regular]\ \ #1}}

\faq{``Generate options'' is disabled}
The period has input errors. Open \ui{Prepare}: the input check lists each error, and a red dot marks the
section that contains it (orders, opening stock or policy). Fix it and click \ui{Save \& check}. Typical causes
are opening stock above a store's capacity, a pull-ahead quantity beyond the after-period demand, or a stage whose
load exceeds its capacity (add calendar time, remove orders, or accept lateness by moving due dates).

\faq{All options have several late orders}
Look at the stage load: a stage above about 85\% over the period (heat treatment in the example) cannot deliver
every order on time. Use the order view (Section~\ref{sec:why-late}) to see where late orders wait. Options:
add capacity for the period (remove a maintenance exception, work a rest day via the calendar), accept pull-ahead
from the previous period, negotiate due dates, or try a longer search (\ui{Thorough}).

\faq{``No improvement found, EDD schedule kept''}
A CP-SAT method ran out of time before improving its starting schedule. Generate again with \ui{Balanced} or
\ui{Thorough}, or use the Hybrid method, which brings late orders down with a fast search before CP-SAT starts.

\faq{The options shown were generated from older inputs}
You saved changes under Prepare after the options were generated. They are kept for reference but marked as
based on the old inputs. Click \ui{Generate options for the new inputs}.

\faq{``The order book changed since this period took its orders''}
Orders were added, changed or cancelled in the book after the period was created. Click \ui{Update the period}
under Prepare to take the changes in (Section~\ref{sec:sync}).

\faq{A new period takes orders that were already delivered}
The previous period was committed but not closed, so its delivered orders are still open in the book. Delete the
new period, open the previous one, go to Execute and use \ui{Close \& start the next period}.

\faq{``The period was changed elsewhere --- saving will conflict''}
Another user (or another tab) saved the same period while you were editing. Discard your changes and redo them
on the new version. PlanWise never overwrites somebody else's save silently.

\faq{The edit cannot be saved}
The independent validator found a problem (warehouse cap, missing semi-finished stock, run beyond the period\dots).
The message names it; click \ui{Show} to select the run concerned, then undo or move the upstream run earlier
(Section~\ref{sec:edit-runs}).

\faq{``Every run is kept: unlock some runs or shrink the frozen zone''}
\ui{Optimise the rest} has nothing it may change. Untick \ui{Keep the runs I edited}, unlock runs, or choose a
smaller frozen zone.

\faq{I cannot edit or commit anything}
Either the period is closed (closed periods are read-only), or you are signed in with a reader key. Ask for an
operator key if you need to make changes.

\faq{``API unreachable'' in the sidebar}
The web page cannot reach the backend. Check that the backend is running and that the address configured for
the frontend is correct, then reload the page. Nothing you saved is lost; unsaved drafts of the factory set-up
are kept in the browser.

\faq{Imported dates are wrong}
Use day-first dates (\code{dd/mm/yyyy}) or ISO dates (\code{yyyy-mm-dd}). American month-first dates such as
12/03/2026 are read as 12 March. Check the \ui{Due} column of the import preview before writing.

\faq{Why are the figures different each time I generate?}
Randomised algorithms (simulated annealing, genetic algorithm, CP-SAT searches) depend on the seed and on how
much the computer can do within the time budget, so two runs can give different, equally valid schedules. The
dispatching rules are deterministic. Every run stores its data snapshot and seed, so a committed schedule can
always be reproduced exactly.
